\documentclass[11pt,a4paper]{article}
\usepackage[T1]{fontenc}
\usepackage{microtype}
\usepackage{amsmath,amssymb,amsthm,mathtools}
\usepackage{mathrsfs}
\usepackage{aliascnt}
\usepackage{booktabs,array,longtable}
\usepackage{geometry}
\usepackage{xcolor}
\usepackage{hyperref}
\usepackage[nameinlink,noabbrev]{cleveref}

\geometry{top=2.45cm,bottom=2.45cm,left=2.75cm,right=2.75cm}
\setlength{\emergencystretch}{2em}
\raggedbottom
\hypersetup{
  colorlinks=true,
  linkcolor=blue!55!black,
  citecolor=blue!55!black,
  urlcolor=blue!55!black,
  pdfauthor={Akari Hayami}
}

\theoremstyle{plain}
\newtheorem{theorem}{Theorem}[section]
\newaliascnt{proposition}{theorem}
\newtheorem{proposition}[proposition]{Proposition}
\aliascntresetthe{proposition}
\newaliascnt{lemma}{theorem}
\newtheorem{lemma}[lemma]{Lemma}
\aliascntresetthe{lemma}
\newaliascnt{corollary}{theorem}
\newtheorem{corollary}[corollary]{Corollary}
\aliascntresetthe{corollary}
\theoremstyle{definition}
\newaliascnt{definition}{theorem}
\newtheorem{definition}[definition]{Definition}
\aliascntresetthe{definition}
\newaliascnt{example}{theorem}
\newtheorem{example}[example]{Example}
\aliascntresetthe{example}
\theoremstyle{remark}
\newaliascnt{remark}{theorem}
\newtheorem{remark}[remark]{Remark}
\aliascntresetthe{remark}
\crefname{proposition}{proposition}{propositions}
\Crefname{proposition}{Proposition}{Propositions}
\crefname{lemma}{lemma}{lemmas}
\Crefname{lemma}{Lemma}{Lemmas}
\crefname{corollary}{corollary}{corollaries}
\Crefname{corollary}{Corollary}{Corollaries}
\crefname{definition}{definition}{definitions}
\Crefname{definition}{Definition}{Definitions}
\crefname{example}{example}{examples}
\Crefname{example}{Example}{Examples}
\crefname{remark}{remark}{remarks}
\Crefname{remark}{Remark}{Remarks}
\numberwithin{equation}{section}

\newcommand{\N}{\mathbb N}
\newcommand{\R}{\mathbb R}
\newcommand{\Z}{\mathbb Z}
\newcommand{\C}{\mathbb C}
\newcommand{\T}{\mathbb T}
\newcommand{\F}{\mathbb F}
\newcommand{\kk}{\Bbbk}
\newcommand{\norm}[1]{\left\lVert#1\right\rVert}
\newcommand{\Aex}{\mathcal A_{\mathrm{ex}}}
\newcommand{\Hex}{H_{\mathrm{ex}}}
\newcommand{\Nex}{\mathcal N_{\mathrm{ex}}}
\DeclareMathOperator{\im}{im}
\DeclareMathOperator{\id}{id}
\DeclareMathOperator{\dist}{dist}
\DeclareMathOperator{\reach}{reach}
\DeclareMathOperator{\Fill}{Fill}
\DeclareMathOperator{\Fillw}{\widehat{Fill}}
\DeclareMathOperator{\End}{End}
\DeclareMathOperator{\supp}{supp}
\DeclareMathOperator{\rank}{rank}
\DeclareMathOperator{\ob}{ob}
\usepackage{graphicx}
\usepackage{tikz}
\usepackage{pgfplots}
\usepackage{caption}
\usepackage[section]{placeins}
\usetikzlibrary{arrows.meta,positioning,calc,fit,backgrounds}
\pgfplotsset{compat=1.18}
\definecolor{viewblue}{HTML}{245A81}
\definecolor{viewteal}{HTML}{16766D}
\definecolor{vieworange}{HTML}{AE591E}
\definecolor{viewgray}{HTML}{53616C}
\tikzset{
  >=Latex,
  every picture/.style={font=\small},
  vbox/.style={draw=viewblue,rounded corners=2pt,fill=viewblue!5,
    align=center,inner sep=7pt,line width=.7pt},
  vinput/.style={vbox,draw=vieworange,fill=vieworange!6},
  vresult/.style={vbox,draw=viewteal,fill=viewteal!6},
  varrow/.style={->,line width=.8pt,draw=viewblue},
  vconditional/.style={->,dashed,line width=.8pt,draw=vieworange},
  vlabel/.style={font=\footnotesize,align=center,fill=white,inner sep=2pt},
  vpoint/.style={circle,fill=viewblue,inner sep=2.5pt},
  vseed/.style={circle,draw=vieworange,fill=vieworange!20,inner sep=3pt}
}
\pgfplotsset{viewaxis/.style={
  axis lines=left,axis line style={viewgray},
  tick label style={font=\scriptsize},label style={font=\small},
  legend style={font=\scriptsize,draw=none,fill=white},
  grid=major,grid style={viewgray!12},
  every axis plot/.append style={line width=.85pt},
  scaled ticks=false
}}
\captionsetup[figure]{font=small,labelfont=bf,labelsep=period}
\renewcommand{\topfraction}{.9}
\renewcommand{\bottomfraction}{.85}
\renewcommand{\textfraction}{.08}
\renewcommand{\floatpagefraction}{.75}
\setcounter{topnumber}{3}
\setcounter{bottomnumber}{2}
\setcounter{totalnumber}{4}
\newcommand{\ViewHeading}[1]{\par\smallskip{\small\bfseries #1}\par\smallskip}
\usepgfplotslibrary{groupplots}
\pgfplotstableread[row sep=\\]{
q boundary cost boundary_threshold cost_threshold \\
1 1.86406484763 0.53646202345 6 0.166666666667 \\
2 1.35098058852 0.740203085444 2.61165168989 0.382899451666 \\
3 0.884941963093 1.1300176076 1.60548312352 0.622865469808 \\
4 1.99234208173 0.501921838208 1.13678742488 0.879671940515 \\
5 0.559007522646 1.78888469204 0.869735596413 1.14977471788 \\
6 1.58720150258 0.63003972613 0.698827118772 1.43096908111 \\
7 1.70933146271 0.58502404116 0.580809354343 1.72173535519 \\
8 0.348363900759 2.87056149567 0.49481546654 2.02095542201 \\
9 1.96180812321 0.509733846125 0.429596009985 2.32776836087 \\
10 1.07345610525 0.931570462087 0.378574406688 2.6414886541 \\
11 1.18382104691 0.844722268293 0.337660320373 2.9615561547 \\
12 1.93143014875 0.517751056463 0.304182170947 3.28750365903 \\
13 0.215982528793 4.63000412853 0.276324394865 3.61893491339 \\
14 1.77489685849 0.563413020433 0.252811955296 3.9555091405 \\
15 1.50233863679 0.665628890525 0.232724303661 4.29692981897 \\
16 0.686077261748 1.4575617875 0.215380941562 4.64293633758 \\
17 1.99957298142 0.500106777444 0.200268453766 4.99329765219 \\
18 0.76311281695 1.31042223088 0.186992524241 5.3478073739 \\
19 1.44650713509 0.69132047519 0.175245521946 5.70627990317 \\
20 1.81146854511 0.552038291086 0.164784081496 6.06854734342 \\
21 0.133645705257 7.48247014804 0.155413269397 6.43445700538 \\
22 1.90832824062 0.524018865684 0.146975191052 6.80386936628 \\
23 1.2494147676 0.800374724174 0.139340654829 7.17665638379 \\
24 1.00281516246 0.997192740427 0.132402980131 7.55270009036 \\
25 1.97620483614 0.506020419397 0.126073334611 7.93189141132 \\
26 0.429438858083 2.32862020094 0.120277178801 8.31412916371 \\
27 1.66496911815 0.60061174054 0.114951523994 8.69931920219 \\
28 1.63612486044 0.611200296615 0.110042795045 9.08737368574 \\
29 0.479187985235 2.08686367524 0.105505148453 9.47821044436 \\
30 1.98341624739 0.504180603196 0.101299136822 9.87175242917 \\
31 0.958292504032 1.04352271962 0.0973906395481 10.2679272324 \\
32 1.28889396577 0.77585901289 0.09375 10.6666666667 \\
33 1.89241812237 0.528424447102 0.0903513243069 11.067906394 \\
34 0.0826356523948 12.1013142756 0.0871719076184 11.4715855982 \\
35 1.83252795147 0.54569426851 0.0841917617004 11.8776466937 \\
36 1.41075988236 0.708837848669 0.0813932236551 12.2860350665 \\
37 0.810080036867 1.23444592446 0.0787606300247 12.6966988416 \\
38 1.99786521099 0.500534267526 0.0762800439225 13.1095886759 \\
39 0.63787233559 1.56771181976 0.0739390254289 13.5246575702 \\
40 1.535567376 0.651225088282 0.0717264374843 13.9418607012 \\
41 1.75077445377 0.571175800428 0.0696322810612 14.3611552682 \\
42 0.266693975521 3.74961600857 0.0676475546085 14.7825003548 \\
43 1.9440608539 0.51438719009 0.0657641337134 15.2058568027 \\
44 1.14226370096 0.875454590002 0.0639746676803 15.6311870973 \\
45 1.11620548774 0.89589238808 0.0622724903268 16.0584552625 \\
46 1.95123338282 0.512496356821 0.0606515427757 16.4876267649 \\
47 0.297950477498 3.35626245139 0.0591063064096 16.9186684255 \\
48 1.73529383975 0.576271278727 0.0576317444676 17.3515483392 \\
49 1.55560430637 0.642836996468 0.0562232510154 17.7862357999 \\
50 0.60786996069 1.64508869441 0.0548766062313 18.2227012324 \\
51 1.99615792675 0.500962367055 0.0535879371159 18.6609161281 \\
52 0.838845057887 1.19211526682 0.0523536828785 19.1008529872 \\
53 1.38820515546 0.720354621988 0.0511705643637 19.5424852634 \\
54 1.84494643912 0.542021155083 0.0500355569822 19.9857873143 \\
55 0.0510806292928 19.5768927252 0.0489458666875 20.4307343536 \\
56 1.88196711397 0.53135891301 0.0478989086067 20.8773024081 \\
57 1.31287462283 0.761687355831 0.0468922879928 21.3254682765 \\
58 0.930461519698 1.07473547141 0.0459237832081 21.7752094915 \\
59 1.98722640468 0.503213925522 0.0449913304941 22.2265042847 \\
60 0.509780380222 1.96162904418 0.0440930103111 22.6793315527 \\
61 1.61776318618 0.618137443441 0.043227035065 23.1336708265 \\
62 1.68225389376 0.594440591702 0.0423917380592 23.5895022423 \\
63 0.398550046504 2.50909517831 0.0415855635313 24.0468065137 \\
64 1.97110295235 0.507330172079 0.0408070576545 24.5055649066 \\
65 1.03000638488 0.970867768079 0.0400548603953 24.9657592145 \\
66 1.22460605134 0.816589138121 0.039327698135 25.427371736 \\
67 1.91753935138 0.521501683542 0.0386243769731 25.8903852533 \\
68 0.165130172021 6.05582848828 0.0379437766404 26.3547830116 \\
69 1.7978612955 0.556216434772 0.037284844958 26.8205487008 \\
70 1.4681297507 0.681138706933 0.0366465927866 27.2876664367 \\
71 0.733834581439 1.36270492737 0.0360280894158 27.756120744 \\
72 1.99997620227 0.500005949503 0.0354284583507 28.2258965406 \\
73 0.715647731484 1.39733552697 0.0348468734549 28.6969791219 \\
74 1.4813106662 0.675077836685 0.0342825554168 29.1693541465 \\
75 1.78922731117 0.558900478301 0.0337347685074 29.6430076222 \\
76 0.184568565779 5.41804069278 0.0332028176025 30.1179258933 \\
77 1.92299333864 0.520022602214 0.0326860454439 30.5940956277 \\
78 1.20912043365 0.827047473659 0.0321838301183 31.0715038056 \\
79 1.04668357016 0.955398583207 0.0316955827342 31.550137708 \\
80 1.96770412102 0.508206487609 0.0312207452776 32.0299849062 \\
81 0.379409558971 2.63567423739 0.0307587886325 32.5110332513 \\
82 1.69272712657 0.590762671847 0.0303092107507 32.9932708649 \\
83 1.60621317182 0.622582367983 0.029871534958 33.4766861297 \\
84 0.528624483397 1.8917020142 0.029445308385 33.9612676805 \\
85 1.98933364995 0.502680885141 0.0290301005132 34.4470043962 \\
86 0.913144642074 1.09511675799 0.0286255018245 34.9338853912 \\
87 1.32753227136 0.753277356473 0.0282311225476 35.4219000082 \\
88 1.8752733639 0.533255587827 0.0278465914929 35.9110378107 \\
89 0.031571686869 31.6739490085 0.0274715549675 36.4012885758 \\
90 1.85239179078 0.539842599702 0.0271056757659 36.8926422877 \\
91 1.37409223001 0.727753187277 0.0267486322285 37.3850891312 \\
92 0.856518749159 1.16751676596 0.0264001173641 37.8786194853 \\
93 1.99485400516 0.501289817406 0.0260598380306 38.3732239175 \\
94 0.589251252537 1.69706894248 0.0257275141696 38.8688931782 \\
95 1.56779428943 0.637838782003 0.0254028780914 39.365618195 \\
96 1.72550981914 0.579538863766 0.025085673805 39.8633900677 \\
97 0.317231427343 3.15227280089 0.0247756563923 40.362200063 \\
98 1.95542322636 0.511398241832 0.0244725914209 40.86203961 \\
99 1.09996112609 0.909123037428 0.0241762543937 41.3629002952 \\
100 1.15822664205 0.863388877182 0.0238864302332 41.8647738585 \\
101 1.93938563029 0.515627209144 0.0236029127966 42.3676521884 \\
102 0.24734266715 4.04297411167 0.0233255044211 42.8715273182 \\
103 1.76012408665 0.568141762042 0.0230540154959 43.3763914221 \\
104 1.522992212 0.656602175716 0.0227882640603 43.8822368108 \\
105 0.65633581581 1.52361028594 0.0225280754249 44.3890559286 \\
106 1.99867145022 0.500332358222 0.0222732818155 44.8968413493 \\
107 0.792200878353 1.26230609852 0.0220237220372 45.405585773 \\
108 1.42452403779 0.701988856256 0.0217792411579 45.9152820225 \\
109 1.82462436158 0.548058011861 0.0215396902098 46.4259230407 \\
110 0.102127932864 9.79164046467 0.0213049259079 46.9375018868 \\
111 1.89864155847 0.526692358302 0.0210748103838 47.450011734 \\
112 1.27391211924 0.784983504667 0.0208492109344 47.9634458659 \\
113 0.975374030104 1.02524771948 0.0206279997848 48.4777976747 \\
114 1.98081423313 0.504842899083 0.0204110538632 48.9930606574 \\
115 0.460220649899 2.17287077453 0.0201982545887 49.5092284142 \\
116 1.64726941829 0.60706523711 0.0199894876703 50.0262946452 \\
117 1.65407879867 0.6045661191 0.0197846429159 50.5442531487 \\
118 0.448476172586 2.22977286448 0.0195836140525 51.0630978183 \\
119 1.97911180616 0.505277163668 0.0193862985546 51.5828226407 \\
120 0.985884673844 1.01431742123 0.0191925974819 52.1034216939 \\
}\ArithmeticProfileData

\pgfplotstableread[row sep=\\]{
q boundary cost \\
1 1.86406484763 0.53646202345 \\
2 1.35098058852 0.740203085444 \\
3 0.884941963093 1.1300176076 \\
5 0.559007522646 1.78888469204 \\
8 0.348363900759 2.87056149567 \\
13 0.215982528793 4.63000412853 \\
21 0.133645705257 7.48247014804 \\
34 0.0826356523948 12.1013142756 \\
}\ArithmeticHitData

\pgfplotstableread[row sep=\\]{
q component first_component \\
1 1.45793725484 1.45793725484 \\
2 0.998026728428 1.45793725484 \\
3 0.424949326499 1.45793725484 \\
4 0.0626666167822 1.45793725484 \\
5 0.32360679775 1.45793725484 \\
6 0.327429083576 1.45793725484 \\
7 0.153093369994 1.45793725484 \\
8 0.0621724717912 1.45793725484 \\
9 0.194734817788 1.45793725484 \\
10 0.190211303259 1.45793725484 \\
11 0.0774144166482 1.45793725484 \\
12 0.0613540921141 1.45793725484 \\
13 0.143042536291 1.45793725484 \\
14 0.129261007495 1.45793725484 \\
15 0.0412022659167 1.45793725484 \\
16 0.0602192092627 1.45793725484 \\
17 0.113950960133 1.45793725484 \\
18 0.0938142139447 1.45793725484 \\
19 0.0197243489038 1.45793725484 \\
20 0.0587785252292 1.45793725484 \\
21 0.0944871144109 1.45793725484 \\
22 0.0700466584342 1.45793725484 \\
23 0.00546004517646 1.45793725484 \\
24 0.0570455921607 1.45793725484 \\
}\ArithmeticSignatureData
\hypersetup{pdftitle={Arithmetic Local-System Filling Spectra}}

\title{\textbf{Arithmetic Local-System Filling Spectra}\\
\large Marked Chain Diagrams, Relative Nonextension,\\
and Positive-Reach Carriers}
\author{Akari Hayami}
\date{Reorganized research manuscript --- 16 September 2026}

\begin{document}
\maketitle

\begin{abstract}
We assemble an explicit marked local-system model in which minimum filling
costs detect simultaneous Diophantine recurrence.  At scale $q$, the
holonomy differences $v_q(x)=(e^{2\pi\mathrm i qx_j}-1)_{j=1}^d$ give
both a chain boundary row and a relative connecting cocycle, with the
chain/cochain coefficient convention stated explicitly.  Their roles are
different: a nonzero row makes the marked chain fillable while making the
marked vertex section nonextendable.  The minimum Euclidean filling cost
is $1/\norm{v_q(x)}_2$, with value $+\infty$ when the row vanishes.

For every $\alpha>1/d$ and $C>0$, the large-cost recurrence locus has
Hausdorff dimension $(d+1)/(1+\alpha)$.  This is a transfer of the classical
simultaneous approximation theorem, not a new Diophantine dimension
formula.  Reverse-divisibility maps are nonexpansive after weighting the
chain norm by $q$.  The normalized boundary signature has exactly the
chord metric of the unit product torus: its first component already
determines that metric.  Product-radius carriers have reach equal to their
smallest radius, and their sub-reach offsets preserve the relative
obstruction for each fixed character.  The selected recurrence locus is
dense and null, with full box dimension.  None of these conclusions
identifies arithmetic scale with offset radius or applies to an unmarked
filtered complex without additional data.
\end{abstract}

\tableofcontents
\clearpage

\section{The common model and the main result}
\label{sec:overview}

The companion \emph{Foundations of Information Exclusion and Relative
Obstruction} separates boundary writing, persistent homology, and extension
of local sections.  The present paper supplies one explicit quantitative
model, rather than a universal passage between those notions.  The
companion \emph{Minimal Marked Response Modules for Relative Obstruction
Data} begins only after a finite quotient, seed, and endomorphisms have
been supplied; those endomorphisms are not outputs of this paper.

The construction has two branches with common coefficients:
\[
\begin{array}{ccc}
 & v_q(x)=(e^{2\pi\mathrm i qx_j}-1)_j & \\
 \swarrow & & \searrow \\
 \partial_{q,x}:\C^d\to\C && d^0_{qx}:\C\to\C^d \\
 \text{minimum chain filling cost} && \text{relative connecting class}.
\end{array}
\]
The arithmetic spectrum comes from the size of $v_q(x)$.  Geometric offset
transport is a further construction for a fixed local system, not a
consequence of the word ``filling''.

Fix an integer $d\geq1$.  Let $\T^d=(\R/\Z)^d$, with metric
\[
 d_{\T}(x,y)=\left(\sum_{j=1}^d
                  \dist(x_j-y_j,\Z)^2\right)^{1/2}.
\]
Distances in this expression are independent of the chosen real lifts.
Write
\begin{equation}\label{eq:basic-data}
 z_{q,j}(x)=e^{2\pi\mathrm i qx_j},\qquad
 v_q(x)=(z_{q,j}(x)-1)_{j=1}^d,\qquad
 b_q(x)=\norm{v_q(x)}_2.
\end{equation}
For $\alpha>0$ and $C>0$, define the recurrence set
\begin{equation}\label{eq:recurrence-locus}
 R_{\alpha,C}
 =\{x\in\T^d:b_q(x)<Cq^{-\alpha}
                       \text{ for infinitely many }q\in\N\}.
\end{equation}
Here $\N=\{1,2,\ldots\}$; denominators need not be reduced.  Our dimension
and nullity results concern $\alpha>1/d$.  The broader definition also
allows the all-positive-exponent rational membership and density statement
in \cref{prop:legacy-density}; it does not extend the dimension theorem.

\begin{theorem}[Marked filling and geometric realization]
\label{thm:main}
For $d\geq1$, $\alpha>1/d$, and $C>0$, the following statements hold.
\begin{enumerate}
 \item There are finite marked two-term local-system chain complexes
 \[
  \mathcal C_x(q):\quad
  \C^d\xrightarrow{\partial_{q,x}}\C,
  \qquad \partial_{q,x}c=\sum_{j=1}^d(z_{q,j}(x)-1)c_j,
 \]
 whose marked zero-chain $1$ has minimum filling cost
 \[
  \Fill_{q,x}(1)=
  \begin{cases}b_q(x)^{-1},&b_q(x)>0,\\+\infty,&b_q(x)=0.\end{cases}
 \]
 \item With degree-one norm $\norm{c}_q=q\norm{c}_2$, the complexes
 form a nonexpansive reverse-divisibility diagram.  The feasible sets
 $\mathcal F_x(q,\lambda)$ form a Set-valued functor, and
 \[
  x\in R_{\alpha,C}
  \iff \mathcal F_x(q,C^{-1}q^{\alpha+1})=\varnothing
                     \text{ for infinitely many }q.
 \]
 \item For $U(x)=(q^{-1}v_q(x))_{q\geq1}$,
 \begin{align*}
  \norm{U(x)-U(y)}_{\ell^\infty(\ell^2)}
    &=\norm{v_1(x)-v_1(y)}_2,\\
  4d_{\T}(x,y)&\leq\norm{U(x)-U(y)}_{\ell^\infty(\ell^2)}
                  \leq2\pi d_{\T}(x,y).
 \end{align*}
 \item The recurrence locus and its signature image have dimension
 \[
  \dim_{\mathrm H}R_{\alpha,C}
   =\dim_{\mathrm H}U(R_{\alpha,C})=\frac{d+1}{1+\alpha}.
 \]
 \item Let $E_{\alpha,C}=R_{\alpha,C}\setminus\{0\}$.  On every
 product-radius carrier $\mathcal M_{\boldsymbol\lambda}$ defined in
 \cref{sec:carriers}, its image is dense and $\mathcal H^d$-null, has
 Hausdorff dimension $(d+1)/(1+\alpha)$, and has both box dimensions
 equal to $d$.  For each $x\in E_{\alpha,C}$, the relative obstruction
 of the original local system $\mathcal L_x$ persists through all
 offsets $0\leq r<\min_j\lambda_j$.
\end{enumerate}
\end{theorem}

The proofs appear in \cref{sec:local,sec:diagram,sec:signature,sec:dimension,sec:carriers}.
The reciprocal formula and signature identity are elementary calculations;
the dimension theorem uses the classical result specified in
\cref{sec:dimension}.  The purpose of the synthesis is to make their common
model and their distinct hypotheses explicit.

\section{Local coefficients, filling, and relative nonextension}
\label{sec:local}

\subsection{Coefficient conventions}

Let $K_d$ be the one-vertex product CW model of $\T^d$, with vertex $a$,
and let $R_d=K_d^{(1)}=\bigvee_{j=1}^d S^1$.  Orient and order the loops,
and identify the vertex fibre with $\C$, retaining its vector $1$.
For $x\in\T^d$, let $\mathcal L_x$ be the rank-one local system with
holonomy
\[
 \rho_x(n)=e^{2\pi\mathrm i\langle x,n\rangle},\qquad n\in\Z^d.
\]
We use the cellular cochain convention
\begin{equation}\label{eq:coboundary}
 d^0_{qx}(s)=v_q(x)s.
\end{equation}
Its restriction to $R_d$ has no cochains above degree one.  Define
$\mathcal C_x(q)$ as the degreewise \emph{algebraic dual} of this wedge
cochain complex.  Equivalently, it is the cellular chain complex with the
dual coefficient system $\mathcal L_{qx}^{\vee}|_{R_d}$, under the
corresponding cellular pairing.  Thus its boundary is the transpose row
in \cref{thm:main}.  This convention distinguishes a local system from its
dual and uses a bilinear evaluation pairing, not a Hermitian adjoint.
Local-coefficient chains, cochains, and relative sequences are discussed in
\cite[Section 3.H]{hatcher}.

The degree-one chain norm is initially Euclidean and the degree-zero norm
is absolute value.  The selected zero-chain $1$ is a cycle because there
is no negative degree.  Put
\begin{equation}\label{eq:filling-definition}
 \Fill_{q,x}(1)=\inf\{\norm{c}_2:\partial_{q,x}c=1\},
 \qquad \inf\varnothing=+\infty.
\end{equation}

\begin{proposition}[Exact filling cost]\label{prop:filling}
The filling formula in \cref{thm:main} holds.  When $b_q(x)>0$, a
minimum-norm filler is
\[
 c_j=\frac{\overline{z_{q,j}(x)-1}}{b_q(x)^2}.
\]
Consequently, for every $\alpha>0$ and $C>0$,
\begin{equation}\label{eq:cost-threshold}
 b_q(x)<Cq^{-\alpha}
 \iff \Fill_{q,x}(1)>C^{-1}q^{\alpha}.
\end{equation}
\end{proposition}

\begin{proof}
The displayed vector has boundary $1$ and norm $1/b_q(x)$.  Every filler
satisfies $1=|\partial_{q,x}c|\leq b_q(x)\norm{c}_2$ by
Cauchy--Schwarz, proving minimality.  If $b_q(x)=0$, the boundary is zero
and no filler exists.  The threshold equivalence includes that case.
\end{proof}

\subsection{The other branch of the same coefficient vector}

Let
\[
 \delta_y:H^0(\{a\};\mathcal L_y)\longrightarrow
                    H^1(K_d,\{a\};\mathcal L_y)
\]
be the connecting homomorphism.  The relative degree-zero cochain space
is zero.  Extending the vertex state $1$ as the absolute zero-cochain $1$
represents $\delta_y(1)$ by $v_1(y)$; it is a cocycle because
$d^1_y d^0_y=0$ and cannot be a relative boundary unless it vanishes.

\begin{proposition}[Fillability and nonextension are different events]
\label{prop:two-events}
For every $x\in\T^d$ and $q\geq1$,
\begin{align}
 \delta_{qx}(1)\ne0
  &\iff qx\ne0\iff b_q(x)>0\iff\Fill_{q,x}(1)<\infty,
  \label{eq:nontrivial-case}\\
 \delta_{qx}(1)=0
  &\iff qx=0\iff b_q(x)=0\iff\Fill_{q,x}(1)=+\infty.
  \label{eq:trivial-case}
\end{align}
In particular, the unpowered vertex state has a nonzero obstruction for
every $x\ne0$, not just for the selected recurrence locus.
\end{proposition}

\begin{proof}
The preceding relative-cochain calculation proves the first equivalence
in each line.  Alternatively, $H^0(K_d;\mathcal L_y)=\ker d^0_y$
vanishes for $y\ne0$, so exactness makes $\delta_y$ injective.  At $y=0$
the vertex state extends constantly.  The remaining equivalences follow
from \cref{eq:basic-data,prop:filling}.
\end{proof}

\Cref{fig:arithmetic-branches} places the two coefficient branches and
their opposite solvability interpretations side by side.
\begin{figure}[tbp]
\centering
\begin{tikzpicture}
\node[vbox,text width=8cm] (vector) at (0,0)
 {$v_q(x)=(e^{2\pi\mathrm i qx_j}-1)_j$,\qquad $b_q(x)=\|v_q(x)\|_2$};
\node[vbox,text width=5.6cm] (chain) at (-3.55,-1.7)
 {\textbf{Chain branch on the wedge}\\
 $\partial_{q,x}:\C^d\to\C$, $c\mapsto\sum_jv_{q,j}c_j$\\
 Algebraic dual coefficients $\mathcal L_{qx}^{\vee}$};
\node[vbox,text width=5.6cm] (cochain) at (3.55,-1.7)
 {\textbf{Cochain branch on the torus}\\
 $d^0_{qx}:\C\to\C^d$, $s\mapsto v_q(x)s$\\Coefficients $\mathcal L_{qx}$};
\draw[varrow] (vector)--(chain);
\draw[varrow] (vector)--(cochain);
\node[vresult,text width=5.6cm] (cost) at (-3.55,-3.55)
 {Minimum norm of a filler of $1$\\$\Fill_{q,x}(1)=b_q(x)^{-1}$ if $b_q(x)>0$};
\node[vresult,text width=5.6cm] (ob) at (3.55,-3.55)
 {Relative connecting class\\$\delta_{qx}(1)=[v_q(x)]$};
\draw[varrow] (chain)--(cost);
\draw[varrow] (cochain)--(ob);
\end{tikzpicture}
\medskip
\begin{tabular}{@{}cccl@{}}
\toprule
Character & $b_q$ & Chain cost & Vertex section\\\midrule
$qx\neq0$ & $>0$ & finite & $\delta_{qx}(1)\neq0$: does not extend\\
$qx=0$ & $0$ & $+\infty$ & $\delta_{qx}(1)=0$: extends\\\bottomrule
\end{tabular}
\caption{Common coefficients, different events. The chain boundary is
the algebraic transpose of the wedge coboundary, not its Hermitian
adjoint. Conjugation in the minimum-norm filler comes from Euclidean
minimization. Non-fillability is not relative nonextension.}
\label{fig:arithmetic-branches}
\end{figure}

\begin{remark}[Three independent roles]\label{rem:three-scales}
The parameter $x$ selects a local system; the integer $q$ replaces it by
its tensor power $\mathcal L_{qx}$; the offset radius $r$ will change the
geometric domain while pulling back a \emph{fixed} local system.  If
$x\ne0$ is torsion, then $\delta_x(1)\ne0$, but $\delta_{qx}(1)=0$
at multiples of its order.  Preservation along sub-reach offsets must not
be described as preservation under every tensor power.
\end{remark}

\begin{example}[Two-dimensional cochains and a powered torsion point]
\label{ex:two-dimensional}
Set $d=2$, $u=e^{2\pi\mathrm i qx_1}$, and
$v=e^{2\pi\mathrm i qx_2}$.  With the product-cell orientation, the full
torus cochain complex is
\[
 \C\xrightarrow{d^0}\C^2\xrightarrow{d^1}\C,
 \qquad
 d^0(s)=((u-1)s,(v-1)s),\quad
 d^1(a,b)=(1-v)a+(u-1)b.
\]
The identity $d^1d^0=0$ is
$(1-v)(u-1)+(u-1)(v-1)=0$.  Relative to the vertex, the degree-zero
space vanishes, so the connecting class of $1$ is represented by
$(u-1,v-1)$ in $\ker d^1$, with no relative boundaries to divide out.
When $(u,v)\ne(1,1)$, the nonzero row $d^1$ has one-dimensional kernel,
and this vector spans it.  At $(u,v)=(1,1)$, the relative $H^1$ is
two dimensional, but the connecting class is zero.

In contrast, the wedge has no degree-two cochains.  Its algebraic dual
has boundary $\partial(a,b)=(u-1)a+(v-1)b$, not the row $d^1$.
This distinction is essential even though both calculations use the same
holonomy differences.

For an explicit specialization take $x=(1/2,0)$.  At $q=1$ the boundary
is $(-2,0)$, the minimum filler is $(-1/2,0)$ with cost $1/2$, and
$\delta_x(1)$ is nonzero.  At $q=2$ both holonomies are $1$, no filler
of the marked zero-chain exists, and $\delta_{2x}(1)=0$.  The original
class $\delta_x(1)$ has not changed: the second calculation uses a
different, powered local system.
\end{example}

\Cref{fig:arithmetic-wedge,fig:arithmetic-torsion} separate the full-torus
cochain calculation from the wedge and display the exact torsion example.
\begin{figure}[tbp]
\centering
\begin{minipage}[t]{.46\linewidth}
\centering\ViewHeading{(a) Wedge $R_2=S^1\vee S^1$}
\begin{tikzpicture}
\draw[viewblue,thick] (-.95,0) circle (.95);
\draw[viewblue,thick] (.95,0) circle (.95);
\node[vseed,label=below:$a$] at (0,0) {};
\draw[varrow] (-1.9,0) arc[start angle=180,end angle=120,radius=.95];
\draw[varrow] (1.9,0) arc[start angle=0,end angle=60,radius=.95];
\node at (-.95,0) {$e_1$};\node at (.95,0) {$e_2$};
\node[align=center] at (0,-1.9)
 {$\C\xrightarrow{d^0}\C^2\longrightarrow0$\\Relative $Z^1=\C^2$};
\end{tikzpicture}
\end{minipage}\hfill
\begin{minipage}[t]{.50\linewidth}
\centering\ViewHeading{(b) Product CW torus $K_2$}
\begin{tikzpicture}
\fill[viewteal!10] (-1,-1) rectangle (1,1);
\draw[varrow] (-1,-1)--(1,-1) node[midway,below] {$e_1$};
\draw[varrow] (-1,1)--(1,1) node[midway,above] {$e_1$};
\draw[varrow] (-1,-1)--(-1,1) node[midway,left] {$e_2$};
\draw[varrow] (1,-1)--(1,1) node[midway,right] {$e_2$};
\foreach \point in {(-1,-1),(-1,1),(1,-1),(1,1)} \fill[vieworange] \point circle (2.5pt);
\node at (0,0) {one $2$-cell};
\node[align=center] at (0,-2)
 {$\C\xrightarrow{d^0}\C^2\xrightarrow{d^1}\C$\\Relative $Z^1=\ker d^1$};
\end{tikzpicture}
\end{minipage}
\[
d^0(s)=((u-1)s,(v-1)s),\qquad d^1(a,b)=(1-v)a+(u-1)b.
\]
\caption{The same degree-zero coboundary does not give the same relative
degree-one quotient. Opposite square edges are identified with the
displayed orientations, and all corners represent one vertex. For a
nontrivial character the torus kernel is one dimensional; the wedge has
no degree-two equation.}
\label{fig:arithmetic-wedge}
\end{figure}
\begin{figure}[tbp]
\centering
\begin{tikzpicture}[x=1.65cm]
\draw[->,viewgray] (.5,0)--(6.6,0) node[right] {$q$};
\foreach \scale in {1,3,5} {
 \node[vseed] at (\scale,0) {};
 \node[below] at (\scale,-.2) {$\scale$};
 \node[vieworange] at (\scale,.65) {$(-2,0)$};
}
\foreach \scale in {2,4,6} {
 \node[vpoint] at (\scale,0) {};
 \node[below] at (\scale,-.2) {$\scale$};
 \node[viewblue] at (\scale,.65) {$(0,0)$};
}
\end{tikzpicture}
\medskip
\begin{tabular}{@{}llll@{}}
\toprule
$x=(1/2,0)$ & $q$ odd & $q$ even & What is held fixed?\\\midrule
$\Fill_{q,x}(1)$ & $1/2$ & $+\infty$ & parameter $x$\\
$\delta_{qx}(1)$ & nonzero & zero & not the powered system\\
$\delta_x(1)$ & nonzero & nonzero & original system $\mathcal L_x$\\\bottomrule
\end{tabular}
\caption{Exact torsion behavior, not a numerical approximation. Labels
above the axis are boundary rows. Tensor powers alternate between a
nontrivial and a trivial character, while the unpowered connecting class
is unchanged.}
\label{fig:arithmetic-torsion}
\end{figure}

\section{The nonexpansive arithmetic diagram}
\label{sec:diagram}

For $k\geq1$, put $S_k(z)=1+z+\cdots+z^{k-1}$.  Wrapping each oriented
loop of $R_d$ around itself $k$ times gives a cellular map $f_k$ with
$f_k^*\mathcal L_{qx}=\mathcal L_{kqx}$ under the marked fibre
identifications.  The degree-one cochain pullback is diagonal with entries
$S_k(z_{q,j})$; algebraic dualization reverses its direction.  Define
\begin{equation}\label{eq:arithmetic-map}
 T_{q,k}:\mathcal C_x(kq)_1\longrightarrow\mathcal C_x(q)_1,
 \qquad (T_{q,k}c)_j=S_k(z_{q,j}(x))c_j.
\end{equation}
The degree-zero map is the identity.

\begin{lemma}[Arithmetic identities]\label{lem:arithmetic}
For all positive integers $q,k,\ell$,
\[
 \partial_{q,x}T_{q,k}=\partial_{kq,x},\qquad
 T_{q,k}T_{kq,\ell}=T_{q,k\ell},\qquad
 \norm{T_{q,k}}_{2\to2}\leq k.
\]
\end{lemma}

\begin{proof}
The first two identities follow coordinatewise from
$z^k-1=(z-1)S_k(z)$ and $S_k(z)S_\ell(z^k)=S_{k\ell}(z)$.
For $|z|=1$, the triangle inequality gives $|S_k(z)|\leq k$.
\end{proof}

Give the chain space at scale $q$ the norm $\norm{c}_q=q\norm{c}_2$.
A marked normed chain complex here includes ordered bases, norms, and a
distinguished zero-cycle.  Its morphisms are degreewise nonexpansive chain
maps carrying that cycle to the distinguished cycle; they need not carry
each basis vector to a basis vector.  The bases remain part of the
coordinate data, not a claimed invariant of unmarked complexes.

\begin{proposition}[Reverse-divisibility functor]\label{prop:diagram}
The complexes $\mathcal C_x(q)$ and maps $(T_{q,k},\id)$ form a functor
from $\N_{\mid}^{\mathrm{op}}$, with arrows $kq\to q$, to finite marked
normed chain complexes.  Its maps are nonexpansive.
\end{proposition}

\begin{proof}
The chain, composition, and identity laws follow from
\cref{lem:arithmetic} and $S_1=1$.  The required norm estimate is
\[
 \norm{T_{q,k}c}_q\leq qk\norm{c}_2=\norm{c}_{kq}.
\]
The degree-zero identity preserves $1$ and its norm.
\end{proof}

\Cref{fig:arithmetic-diagram} records the arrow direction, chain square,
and scale-dependent norms.
\begin{figure}[tbp]
\centering
\begin{minipage}[t]{.42\linewidth}
\centering\ViewHeading{(a) Reverse divisibility}
\begin{tikzpicture}
\node[vbox] (six) at (0,0) {$6$};
\node[vbox] (two) at (-1.7,-1.5) {$2$};
\node[vbox] (three) at (1.7,-1.5) {$3$};
\node[vresult] (one) at (0,-3) {$1$};
\draw[varrow] (six)--(two) node[midway,above left,font=\footnotesize] {$T_{2,3}$};
\draw[varrow] (six)--(three) node[midway,above right,font=\footnotesize] {$T_{3,2}$};
\draw[varrow] (two)--(one) node[midway,below left,font=\footnotesize] {$T_{1,2}$};
\draw[varrow] (three)--(one) node[midway,below right,font=\footnotesize] {$T_{1,3}$};
\node[font=\footnotesize] at (0,-3.85) {both composites are $T_{1,6}$};
\end{tikzpicture}
\end{minipage}\hfill
\begin{minipage}[t]{.55\linewidth}
\centering\ViewHeading{(b) Boundary and norm compatibility}
\begin{tikzpicture}
\node (sourceone) at (0,0) {$\mathcal C_x(kq)_1$};
\node (sourcezero) at (4.3,0) {$\C$};
\node (targetone) at (0,-1.65) {$\mathcal C_x(q)_1$};
\node (targetzero) at (4.3,-1.65) {$\C$};
\draw[varrow] (sourceone)--(sourcezero) node[midway,above] {$\partial_{kq,x}$};
\draw[varrow] (targetone)--(targetzero) node[midway,below] {$\partial_{q,x}$};
\draw[varrow] (sourceone)--(targetone) node[midway,left] {$T_{q,k}$};
\draw[varrow] (sourcezero)--(targetzero) node[midway,right] {$\id$};
\node[vinput,text width=6.2cm] at (2.05,-3.25)
 {$\|T_{q,k}\|_{2\to2}\leq k$\\$\|T_{q,k}c\|_q\leq qk\|c\|_2=\|c\|_{kq}$};
\end{tikzpicture}
\end{minipage}
\caption{Arithmetic maps point from a multiple to a divisor; they are not
inclusions of a filtration. Weighting the degree-one norm at each node
by its scale makes the maps nonexpansive. The degree-zero map fixes the
marked chain $1$.}
\label{fig:arithmetic-diagram}
\end{figure}

Define the weighted cost and bounded feasible sets by
\begin{align}
 \Fillw_{q,x}(1)&=\inf\{q\norm{c}_2:\partial_{q,x}c=1\}
                   =q\Fill_{q,x}(1),\label{eq:weighted-cost}\\
 \mathcal F_x(q,\lambda)&=\{c:\partial_{q,x}c=1,
                                  q\norm{c}_2\leq\lambda\},
                                  \qquad\lambda\geq0.
 \label{eq:feasible}
\end{align}
Use the product preorder
$\mathsf P=\N_{\mid}^{\mathrm{op}}\times\R_{\geq0}$, where
$(kq,\lambda')\to(q,\lambda)$ when $\lambda'\leq\lambda$.

\begin{proposition}[Feasible-set functor]\label{prop:feasible}
Restriction of $T_{q,k}$ makes $\mathcal F_x$ a functor
$\mathsf P\to\mathbf{Set}$.  Its values are empty or closed convex
sets, and
\[
 \mathcal F_x(q,\lambda)=\varnothing\iff\Fillw_{q,x}(1)>\lambda.
\]
In particular, the feasible-set equivalence in \cref{thm:main} holds.
\end{proposition}

\begin{proof}
If $c\in\mathcal F_x(kq,\lambda')$, the chain identity and
nonexpansiveness give
\[
 \partial_{q,x}T_{q,k}c=1,\qquad
 \norm{T_{q,k}c}_q\leq\lambda'\leq\lambda.
\]
Composition is inherited from the chain maps.  Closedness and convexity
follow by intersecting an affine subspace with a closed norm ball.  A
finite minimum is attained by \cref{prop:filling}, proving the strict
inequality criterion, including equality at the threshold.  Substitute
$\lambda=C^{-1}q^{\alpha+1}$ and use \cref{eq:weighted-cost}.
\end{proof}

The threshold geometry and induced feasible-set maps are illustrated in
\cref{fig:arithmetic-filler}.
\begin{figure}[tbp]
\centering
\ViewHeading{(a) Real slice of the model row $\partial=(1,0)$}
\begin{tikzpicture}[scale=.9]
\foreach \shift/\radius/\heading in {-4.8/.65/{below cost},0/1/{at cost},4.8/1.35/{above cost}} {
 \begin{scope}[xshift=\shift cm]
 \draw[->,viewgray!60] (-1.5,0)--(1.65,0);
 \draw[->,viewgray!60] (0,-1.55)--(0,1.65);
 \filldraw[draw=viewblue,fill=viewblue!5] (0,0) circle (\radius);
 \draw[vieworange,thick] (1,-1.6)--(1,1.65);
 \fill[viewgray] (0,0) circle (1.8pt);
 \node[font=\footnotesize] at (0,2) {\heading};
 \node[font=\footnotesize,rotate=90,above] at (1.08,.3) {$c_1=1$};
 \end{scope}
}
\fill[viewteal] (1,0) circle (3pt);
\draw[viewteal,line width=2pt] (5.8,-.907)--(5.8,.907);
\node at (-4.8,-2.2) {$\mathcal F=\varnothing$};
\node at (0,-2.2) {$\mathcal F=\{(1,0)\}$};
\node at (4.8,-2.2) {nonempty intersection};
\end{tikzpicture}
\ViewHeading{(b) General complex row and weighted threshold}
\[
c_{\min}=\frac{\overline{v_q}}{b_q^2},\quad
\|c_{\min}\|_2=\frac1{b_q},\quad
\mathcal F_x(q,\lambda)=\varnothing
\ \Longleftrightarrow\ \lambda<\frac q{b_q}\qquad(b_q>0).
\]
\begin{tikzpicture}
\node[vbox] (source) at (-3.9,0) {$\mathcal F_x(kq,\lambda')$};
\node[vresult] (target) at (3.9,0) {$\mathcal F_x(q,\lambda)$};
\draw[varrow] (source)--(target) node[midway,above] {$T_{q,k}|_{\mathcal F}$}
 node[midway,below,font=\footnotesize] {$\lambda'\leq\lambda$};
\end{tikzpicture}
\caption{A ball first meets the affine solution space at the minimum-norm
filler. The top panels are real slices for row $(1,0)$, with ball radii
$0.65,1,1.35$, not all complex dimensions. The closed feasible set is
already nonempty at equality. At $b_q=0$ it is empty for every finite
threshold. Feasible sets are sets of chains, not chain subcomplexes.}
\label{fig:arithmetic-filler}
\end{figure}

This is a diagram over a preorder, in the broad setting of generalized
persistence \cite{bubenik-desilva-scott}, not ordinary persistent homology
of an inclusion filtration.  For $d>1$ the wedge power map is not a
covering; the torus multiplication map is.  The feasible sets are not
chain subcomplexes.  Moreover, $q\mapsto(q,C^{-1}q^{\alpha+1})$ does
not define a chain in $\mathsf P$: for $k>1$ its threshold at $kq$ is
larger, whereas the arrow to $q$ requires a nonincreasing threshold.

\section{The signature metric is the product-torus chord metric}
\label{sec:signature}

Define
\[
 U(x)=(q^{-1}v_q(x))_{q\geq1}\in\ell^\infty(\N;\C^d),
 \qquad
 \iota_d(x)=(e^{2\pi\mathrm i x_1},\ldots,e^{2\pi\mathrm i x_d})
                 \in\C^d\cong\R^{2d}.
\]
The signature is bounded because $b_q(x)\leq2\sqrt d$.

\begin{lemma}[Chord comparison]\label{lem:chord}
For $0\leq s\leq1/2$,
\[
 4s\leq2\sin(\pi s)\leq2\pi s.
\]
Thus chord distance on the unit circle lies between four and $2\pi$
times quotient distance on $\R/\Z$.
\end{lemma}

\begin{proof}
The upper bound is $\sin t\leq t$.  Concavity of sine on $[0,\pi/2]$
places its graph above the chord joining its endpoints.
\end{proof}

\begin{theorem}[Exact signature metric]\label{thm:signature}
For all $x,y\in\T^d$,
\begin{align}
 \norm{U(x)-U(y)}_{\ell^\infty(\ell^2)}
 &=\norm{v_1(x)-v_1(y)}_2
   =\norm{\iota_d(x)-\iota_d(y)}_2,\label{eq:signature-isometry}\\
 4d_{\T}(x,y)&\leq\norm{U(x)-U(y)}_{\ell^\infty(\ell^2)}
                 \leq2\pi d_{\T}(x,y).\label{eq:signature-bounds}
\end{align}
Consequently, $U$ is injective, its inverse on its image is
$1/4$-Lipschitz, and evaluation at $q=1$ is an isometry from its image
to the translated unit product torus.
\end{theorem}

\begin{proof}
For $|z|=|w|=1$,
\[
 z^q-w^q=(z-w)\sum_{h=0}^{q-1}z^{q-1-h}w^h,
 \qquad q^{-1}|z^q-w^q|\leq|z-w|.
\]
Apply this coordinatewise and take the Euclidean norm.  Every component
of the signature difference is bounded by its $q=1$ component, which
itself occurs in the supremum.  This proves the first equality.  The
second follows because $v_1(x)=\iota_d(x)-(1,\ldots,1)$.
Apply \cref{lem:chord} coordinatewise to obtain
\cref{eq:signature-bounds}; the final assertions follow.
\end{proof}

The complex-valued marked signature and the scalar profile $(b_q(x))_q$
must be distinguished.  The latter determines filling costs, but it need
not determine $x$: for example $b_q(-x)=b_q(x)$ for every $q$.
The signature metric is not a new independent metric extracted from
infinitely many scales.  The arithmetic question lies in the recurrence
of its components, while its metric is already determined at scale one.
No metric on a quotient by arbitrary changes of marking is asserted.

\Cref{fig:arithmetic-signature} illustrates the first-component metric
identity and the loss of information in the scalar profile.
\begin{figure}[tbp]
\centering
\ViewHeading{(a) The first component already gives the signature metric}
\begin{tikzpicture}
\begin{axis}[viewaxis,width=12.5cm,height=4.6cm,xmin=.5,xmax=24.5,
 xlabel={$q$},ylabel={$q^{-1}|v_q(x)-v_q(y)|$},ymin=0,ymax=1.7,legend pos=north east]
\addplot[viewblue,ycomb,mark=*,mark size=1.4pt] table[x=q,y=component] {\ArithmeticSignatureData};
\addlegendentry{$x=0.17,\ y=0.43$}
\addplot[vieworange,dashed] table[x=q,y=first_component] {\ArithmeticSignatureData};
\addlegendentry{$|v_1(x)-v_1(y)|$}
\end{axis}
\end{tikzpicture}
\ViewHeading{(b) Scalar costs forget information retained by the signature}
\begin{tikzpicture}
\node[vbox] (positive) at (-4,0) {$x$};
\node[vbox] (negative) at (-4,-1.2) {$-x$};
\node[vinput,text width=4.7cm] (scalar) at (1,-.6) {same scalar profile\\$(b_q(x))_q=(b_q(-x))_q$};
\draw[varrow] (positive)--(scalar);
\draw[varrow] (negative)--(scalar);
\node[align=center,font=\footnotesize,text width=3.4cm] at (5.1,-.6)
 {If $x\neq-x$, then\\$U(x)\neq U(-x)$.};
\end{tikzpicture}
\caption{Panel (a) samples the general inequality proved in the text;
the equality of the supremum with the first component holds at every
scale, not just the 24 plotted ones. The scalar profile determines costs
but is not an injective parameterization. The full marked signature is.}
\label{fig:arithmetic-signature}
\end{figure}

\section{Recurrence exponents and Hausdorff dimension}
\label{sec:dimension}

For $\alpha,c>0$ and any lift $\widetilde x\in\R^d$ of $x$, put
\[
 \Delta_q(x)=\min_{p\in\Z^d}\norm{q\widetilde x-p}_\infty,
 \qquad
 W_{\alpha,c}=\{x:\Delta_q(x)<cq^{-\alpha}
                                  \text{ infinitely often}\}.
\]
Both are independent of the lift.  The same integer $q$ is used in every
coordinate.  In a fundamental domain, the latter inequality is
$\norm{\widetilde x-p/q}_\infty<cq^{-(1+\alpha)}$.

\begin{lemma}[Comparison with simultaneous approximation]
\label{lem:recurrence}
For every $x$ and $q$,
\begin{equation}\label{eq:recurrence-comparison}
 4\Delta_q(x)\leq b_q(x)\leq2\pi\sqrt d\,\Delta_q(x).
\end{equation}
Consequently,
\begin{equation}\label{eq:locus-sandwich}
 W_{\alpha,C/(2\pi\sqrt d)}\subseteq R_{\alpha,C}
                                      \subseteq W_{\alpha,C/4}.
\end{equation}
\end{lemma}

\begin{proof}
Choose nearest integers coordinatewise, with discrepancies
$s_j\in[0,1/2]$.  Then $\Delta_q=\norm{s}_\infty$ and
\cref{lem:chord} gives $4\norm{s}_2\leq b_q\leq2\pi\norm{s}_2$.
Use $\norm{s}_\infty\leq\norm{s}_2\leq\sqrt d\norm{s}_\infty$.
\end{proof}

We use the simultaneous Jarn\'ik--Besicovitch theorem in the following
normalization:
\begin{equation}\label{eq:classical-dimension}
 \dim_{\mathrm H}W_{\alpha,1}=\frac{d+1}{1+\alpha},
                       \qquad\alpha>\frac1d.
\end{equation}
For the systems-of-linear-forms formulation, this is the $n=1,m=d$
specialization of Bovey--Dodson \cite{bovey-dodson}; the same normalization
is displayed in \cite[Section 3.3, equation (11)]{beresnevich-velani}.
Passing from a fundamental cube to the torus uses a finite cover by
coordinate charts with isometric lifts, so does not change the dimension.

\begin{lemma}[Positive constants do not change the dimension]
\label{lem:constants}
For $c>0$ and $\alpha>1/d$,
$\dim_{\mathrm H}W_{\alpha,c}=(d+1)/(1+\alpha)$.
\end{lemma}

\begin{proof}
For $0<\varepsilon<\alpha-1/d$, all sufficiently large $q$ satisfy
$q^{-(\alpha+\varepsilon)}<cq^{-\alpha}<q^{-(\alpha-\varepsilon)}$.
Therefore
\[
 W_{\alpha+\varepsilon,1}\subseteq W_{\alpha,c}
                                  \subseteq W_{\alpha-\varepsilon,1}.
\]
Use \cref{eq:classical-dimension} and let $\varepsilon\downarrow0$.
\end{proof}

\begin{theorem}[Dimension transfer on the full stated range]
\label{thm:dimension}
For $\alpha>1/d$ and $C>0$,
\[
 \dim_{\mathrm H}R_{\alpha,C}
 =\dim_{\mathrm H}U(R_{\alpha,C})=\frac{d+1}{1+\alpha}<d.
\]
Both sets have zero $d$-dimensional Hausdorff measure in their respective
metrics.
\end{theorem}

\begin{proof}
Both outer sets in \cref{eq:locus-sandwich} have the displayed dimension
by \cref{lem:constants}; monotonicity gives the dimension of the middle
set.  The signature map is bi-Lipschitz by \cref{thm:signature}, so
preserves Hausdorff dimension.  The strict inequality implies nullity.
\end{proof}

The common-denominator geometry and the two-sided comparison of loci
are summarized in \cref{fig:arithmetic-recurrence}.
\begin{figure}[tbp]
\centering
\begin{minipage}[t]{.40\linewidth}
\centering\ViewHeading{(a) One common denominator}
\begin{tikzpicture}[scale=.8]
\draw[viewgray!50,step=1] (0,0) grid (4,4);
\foreach \horizontal in {0,...,4}
 \foreach \vertical in {0,...,4} {
 \draw[viewblue!50,fill=viewblue!8] (\horizontal-.15,\vertical-.15) rectangle (\horizontal+.15,\vertical+.15);
 \fill[viewblue] (\horizontal,\vertical) circle (1.6pt);
 }
\node[vseed,label=above right:$x$] at (2.12,1.1) {};
\node[font=\footnotesize,align=center] at (2,-.8)
 {$p/q$, schematic $q=4$ grid\\neighborhood radius $cq^{-(1+\alpha)}$};
\end{tikzpicture}
\end{minipage}\hfill
\begin{minipage}[t]{.57\linewidth}
\centering\ViewHeading{(b) Comparison, not equality of loci}
\begin{tikzpicture}
\node[vbox,text width=6.8cm] (bound) at (0,0)
 {$4\Delta_q(x)\leq b_q(x)\leq2\pi\sqrt d\,\Delta_q(x)$};
\node[vresult,text width=6.8cm] (sandwich) at (0,-1.7)
 {$\displaystyle W_{\alpha,C/(2\pi\sqrt d)}\subseteq R_{\alpha,C}\subseteq W_{\alpha,C/4}$};
\draw[varrow] (bound)--(sandwich);
\node[vinput,text width=6.8cm] at (0,-3.55)
 {Classical dimension theorem\\and constant comparison\\
 $\displaystyle\dim_{\rm H}R_{\alpha,C}=\frac{d+1}{1+\alpha}$\\only in the stated range $\alpha>1/d$};
\end{tikzpicture}
\end{minipage}
\caption{The same integer denominator is used in every coordinate.
Panel (a) is a finite grid schematic, not a picture of the limsup set.
The two comparison constants in (b) are essential: equality of the
resulting Hausdorff dimensions does not identify the selected sets at
the same tolerance.}
\label{fig:arithmetic-recurrence}
\end{figure}

\begin{remark}[Exponent dictionary]\label{rem:exponents}
Earlier toric notation used $\tau=\alpha$, whereas the earlier
quantitative manuscript used $\tau=\alpha-1>0$.  The latter restriction
only covered $\alpha>1$.  The present proof treats the larger range
$\alpha>1/d$; this is a range extension, not merely a renaming.
The consistent threshold dictionary is
\[
\begin{array}{c|c}
 \text{quantity} & \text{condition at scale }q\\ \hline
 b_q(x) & <Cq^{-\alpha}\\
 \Fill_{q,x}(1) & >C^{-1}q^\alpha\\
 \Fillw_{q,x}(1) & >C^{-1}q^{\alpha+1}\\
 \norm{q^{-1}v_q(x)}_2 & <Cq^{-(\alpha+1)}.
\end{array}
\]
The $\Delta_q$ condition is comparable with adjusted constants, not
literally identical at the same constant.  Critical and subcritical
exponents $\alpha\leq1/d$ are outside the stated dimension theorem.
\Cref{sec:legacy} specifies the earlier constant-one locus without
identifying it with the chord-norm locus at an unchanged tolerance.
\end{remark}

\subsection{What Boolean exact fillability loses}

\Cref{fig:arithmetic-profiles} plots the reciprocal threshold equivalence
at finitely many integer scales, without testing an infinitely-often
condition.
\begin{figure}[tbp]
\centering
\begin{tikzpicture}
\begin{groupplot}[viewaxis,
 group style={group size=1 by 2,vertical sep=7mm},
 width=11.8cm,height=3.1cm,scale only axis,
 xmin=1,xmax=120,xtick={1,20,40,60,80,100,120},ymode=log]
\nextgroupplot[
 ylabel={$b_q(x)$},xticklabels=\empty,legend pos=south west,ymin=.005,ymax=8]
\addplot[viewblue,mark=*,mark size=.7pt] table[x=q,y=boundary] {\ArithmeticProfileData};
\addlegendentry{$b_q(x)$}
\addplot[vieworange,dashed] table[x=q,y=boundary_threshold] {\ArithmeticProfileData};
\addlegendentry{$Cq^{-\alpha}$}
\addplot[only marks,mark=o,mark size=2.4pt,vieworange,forget plot]
 table[x=q,y=boundary] {\ArithmeticHitData};
\nextgroupplot[
 xlabel={integer scale $q$},ylabel={$\Fill_{q,x}(1)$},legend pos=north west,
 ymin=.12,ymax=80]
\addplot[viewteal,mark=*,mark size=.7pt] table[x=q,y=cost] {\ArithmeticProfileData};
\addlegendentry{$1/b_q(x)$}
\addplot[vieworange,dashed] table[x=q,y=cost_threshold] {\ArithmeticProfileData};
\addlegendentry{$C^{-1}q^\alpha$}
\addplot[only marks,mark=o,mark size=2.4pt,vieworange,forget plot]
 table[x=q,y=cost] {\ArithmeticHitData};
\end{groupplot}
\end{tikzpicture}
\caption{A finite-scale illustration for $d=1$, $x=(\sqrt5-1)/2$,
$\alpha=1.2$, $C=6$, and $1\leq q\leq120$. A point lies below the upper
threshold exactly when its reciprocal lies above the lower threshold.
Orange rings mark these same scales in both panels. Lines join integer
samples only as visual guides. Data use double
precision and contain no zero boundary norms; infinitely-often membership
cannot be inferred from this finite plot.}
\label{fig:arithmetic-profiles}
\end{figure}

Set $\epsilon_q(x)=1$ when $\Fill_{q,x}(1)<\infty$ and $0$ otherwise.
Call $x$ torsion when $qx=0$ for some positive integer $q$.

\begin{proposition}[Boolean no-go on the full exponent range]
\label{prop:boolean}
All non-torsion parameters have $(\epsilon_q(x))_q=(1,1,\ldots)$.
For every $\alpha>1/d$ and $C>0$, membership in $R_{\alpha,C}$ is not
determined by that sequence or by any function of it.
\end{proposition}

\begin{proof}
The first assertion is \cref{eq:trivial-case}.  The torsion parameters
form a countable set.  By \cref{thm:dimension}, $R_{\alpha,C}$ has
strictly positive Hausdorff dimension, so contains a non-torsion point.
It is also $\mathcal H^d$-null, whereas $\T^d$ has positive
$\mathcal H^d$ measure.  Removing both it and the countable torsion set
therefore leaves a non-torsion point outside it.  These two points have
identical Boolean sequences and different recurrence membership.
\end{proof}

This argument uses the full-range dimension theorem, rather than a
one-coordinate badly approximable example whose lower bound only excludes
recurrence at exponents greater than one.  The original and powered
relative-obstruction Boolean tests do not recover the recurrence rate
either: the original test only removes $0$, and the powered tests are
the same nontriviality tests as \cref{eq:nontrivial-case}.

\begin{example}[Explicit parameters with identical Boolean data]
\label{ex:explicit-arithmetic}
Let $\ell=\sum_{n\geq1}10^{-n!}$ and $\gamma=(\sqrt5-1)/2$, and put
\[
 x_{\mathrm L}=(\ell,0,\ldots,0),\qquad
 x_{\mathrm G}=(\gamma,0,\ldots,0)\quad\text{in }\T^d.
\]
Both are non-torsion, so every exact-fillability indicator is $1$.
For every $C>0$, $x_{\mathrm L}\in R_{\alpha,C}$ for all $\alpha>0$,
whereas $x_{\mathrm G}\notin R_{\alpha,C}$ whenever $\alpha>1$.
The estimates and irrationality check are given in \cref{sec:explicit-estimates}.
This constructive comparison illustrates \cref{prop:boolean} on the
smaller range $\alpha>1$; it does not replace its proof for
$1/d<\alpha\leq1$.
\end{example}

\Cref{fig:arithmetic-boolean} distinguishes the full-range Boolean no-go
argument from its more restricted explicit arithmetic comparison.
\begin{figure}[tbp]
\centering
\begin{tikzpicture}
\node[vbox,text width=5.6cm] (inside) at (-3.55,0)
 {Non-torsion $x_{\rm in}\in R_{\alpha,C}$\\exists since $\dim_{\rm H}R_{\alpha,C}>0$};
\node[vbox,text width=5.6cm] (outside) at (3.55,0)
 {Non-torsion $x_{\rm out}\notin R_{\alpha,C}$\\exists since $R_{\alpha,C}$ is null};
\node[vinput,text width=8cm] (boolean) at (0,-1.8)
 {Identical Boolean data: $(\epsilon_q)_q=(1,1,1,\ldots)$};
\draw[varrow] (inside)--(boolean);
\draw[varrow] (outside)--(boolean);
\node[font=\footnotesize,align=center] at (0,-2.8)
 {Different recurrence membership cannot be recovered from this sequence.\\
 Existence argument: the full stated range $\alpha>1/d$.};
\end{tikzpicture}
\medskip
\begin{tabular}{@{}lll@{}}
\toprule
Explicit comparison & Exact fillability & Recurrence conclusion\\\midrule
$x_{\rm L}=(\sum_{n\geq1}10^{-n!},0,\ldots,0)$ & every scale & in $R_{\alpha,C}$ for all $\alpha>0$\\[3pt]
$x_{\rm G}=((\sqrt5-1)/2,0,\ldots,0)$ & every scale & outside $R_{\alpha,C}$ for $\alpha>1$\\\bottomrule
\end{tabular}
\caption{Boolean solvability forgets the size of a filler. The top
argument works throughout the theorem's range; the constructive pair
below separates membership on the smaller range $\alpha>1$, for every
$C>0$. No finite decimal truncation of the Liouville parameter is used.}
\label{fig:arithmetic-boolean}
\end{figure}

\section{Positive-reach carriers and offset transport}
\label{sec:carriers}

Choose radii $\boldsymbol\lambda=(\lambda_1,\ldots,\lambda_d)>0$ and
write $\lambda_- =\min_j\lambda_j$, $\lambda_+=\max_j\lambda_j$.
Define
\[
 \iota_{\boldsymbol\lambda}(x)
       =(\lambda_j e^{2\pi\mathrm i x_j})_{j=1}^d,
 \quad \mathcal M_{\boldsymbol\lambda}=\iota_{\boldsymbol\lambda}(\T^d),
 \quad a_{\boldsymbol\lambda}=\iota_{\boldsymbol\lambda}(0).
\]
The closed offset of radius $r$ is
$\mathcal M_{\boldsymbol\lambda}^r
 =\{y:\dist(y,\mathcal M_{\boldsymbol\lambda})\leq r\}$.
Reach means the supremum of the radii on which every point at smaller
distance has a unique nearest point \cite{federer}.  The general
closed-positive-reach retraction is proved in the geometric appendix of
\emph{Foundations of Information Exclusion and Relative Obstruction}.
Here we additionally compute the exact reach and metric constants for
the specified product carrier.

\begin{proposition}[Exact reach and metric bounds]\label{prop:reach}
The carrier has reach $\lambda_-$, and
\[
 4\lambda_-d_{\T}(x,y)
 \leq\norm{\iota_{\boldsymbol\lambda}(x)-\iota_{\boldsymbol\lambda}(y)}_2
 \leq2\pi\lambda_+d_{\T}(x,y).
\]
For $0\leq r<\lambda_-$, radial nearest-point projection
$p_r:\mathcal M_{\boldsymbol\lambda}^r\to\mathcal M_{\boldsymbol\lambda}$
is a strong deformation retraction fixing $a_{\boldsymbol\lambda}$.
\end{proposition}

\begin{proof}
For $y=(y_1,\ldots,y_d)\in\C^d$, independent minimization gives
\begin{equation}\label{eq:distance}
 \dist(y,\mathcal M_{\boldsymbol\lambda})^2
                  =\sum_{j=1}^d(|y_j|-\lambda_j)^2.
\end{equation}
At distance less than $\lambda_-$ no $y_j$ vanishes, and the unique
nearest point has coordinates $\lambda_j y_j/|y_j|$.  Conversely, set a
block of smallest radius to zero and put every other block on its circle.
The resulting point has distance $\lambda_-$ and a circle of nearest
points.  This proves the exact reach.

The homotopy
\[
 H_t(y)_j=\left((1-t)+\frac{t\lambda_j}{|y_j|}\right)y_j
\]
is defined throughout every sub-reach offset.  It fixes the carrier and
multiplies each radial deviation in \cref{eq:distance} by $1-t$, so stays
in that offset.  The metric bounds follow from \cref{lem:chord}, after
squaring and summing with weights $\lambda_j^2$.
\end{proof}

The exact circle-block geometry and the witness for loss of uniqueness
are shown in \cref{fig:arithmetic-carrier}.
\begin{figure}[tbp]
\centering
\begin{minipage}[t]{.48\linewidth}
\centering\ViewHeading{(a) One circle block and a sub-reach offset}
\begin{tikzpicture}[scale=.95]
\fill[viewblue!9,even odd rule] (0,0) circle (2) (0,0) circle (1);
\draw[viewblue!45] (0,0) circle (2);
\draw[viewblue!45] (0,0) circle (1);
\draw[viewblue,very thick] (0,0) circle (1.5);
\draw[viewgray,dashed] (0,0)--(1.5,0) node[midway,below] {$\lambda$};
\node[vseed,label=above right:$y$] (point) at (50:1.85) {};
\node[vpoint,label=below left:{$p(y)$}] (projection) at (50:1.5) {};
\draw[varrow] (point)--(projection);
\draw[<->,vieworange] (-2,0)--(-1.5,0) node[midway,above] {$r$};
\node[font=\footnotesize,align=center] at (0,-2.6)
 {$0\leq r<\lambda$; radial contraction\\fixes the circle pointwise};
\end{tikzpicture}
\end{minipage}\hfill
\begin{minipage}[t]{.48\linewidth}
\centering\ViewHeading{(b) Failure of uniqueness at the center}
\begin{tikzpicture}[scale=.95]
\draw[viewblue,very thick] (0,0) circle (1.5);
\foreach \angle in {0,45,...,315} \draw[dashed,vieworange] (0,0)--(\angle:1.5);
\node[vseed,label=below:{$y_j=0$}] at (0,0) {};
\node[font=\footnotesize,align=center] at (0,-2.6)
 {Every point on this circle is nearest.\\Choose a block with radius $\lambda_-$.};
\end{tikzpicture}
\end{minipage}
\[
\operatorname{dist}(y,\mathcal M_{\boldsymbol\lambda})^2
 =\sum_j(|y_j|-\lambda_j)^2,\qquad
\operatorname{reach}(\mathcal M_{\boldsymbol\lambda})=\min_j\lambda_j.
\]
\caption{Exact circle-block geometry underlying the product carrier in
$\R^{2d}$. In (a), $\lambda=1.5$ and $r=0.5$. For the witness in (b),
all other blocks lie on their circles. For $d=2$ the carrier is a product
in $\R^4$, not a standard doughnut embedded in $\R^3$.}
\label{fig:arithmetic-carrier}
\end{figure}

Identify $\mathcal L_x$ with its local system on the carrier and set
$\mathcal L_{x,r}=p_r^*\mathcal L_x$.  These are topological local
systems on the offsets; no unprovided cellular discretization is needed.

\begin{proposition}[Constant sub-reach relative system]
\label{prop:offset}
Fix $x\in\T^d$ and $0\leq r\leq R<\lambda_-$.  If
$j_{r,R}:\mathcal M_{\boldsymbol\lambda}^r
\hookrightarrow\mathcal M_{\boldsymbol\lambda}^R$ is inclusion, then
\[
 j_{r,R}^*:H^1(\mathcal M_{\boldsymbol\lambda}^R,
                  \{a_{\boldsymbol\lambda}\};\mathcal L_{x,R})
 \xrightarrow{\ \cong\ }
 H^1(\mathcal M_{\boldsymbol\lambda}^r,
                  \{a_{\boldsymbol\lambda}\};\mathcal L_{x,r})
\]
preserves the connecting class of $1$.  This class is nonzero exactly when
$x\ne0$.  The system is naturally isomorphic to a constant contravariant
system on $[0,\lambda_-)$.
\end{proposition}

\begin{proof}
The radial projections satisfy $p_Rj_{r,R}=p_r$ and are homotopy
equivalences of pairs fixing the marked vertex.  Their pullbacks are
cohomology isomorphisms, and
$j_{r,R}^*p_R^*=p_r^*$ identifies the restriction maps with identities
on carrier cohomology.  Naturality of the relative sequence preserves
the connecting classes.  Apply \cref{prop:two-events} at $q=1$.
\end{proof}

The proof stops at $r<\lambda_-$ because nearest points cease to be
unique at the displayed medial-axis witness.  This is a boundary of the
projection argument, not a claim that every obstruction dies precisely
at $r=\lambda_-$.

\subsection{Dense-null spectra and calibrated codimension}

\Cref{fig:arithmetic-offset} displays the geometric and cohomological
arrow directions for a fixed character, including the strict sub-reach
bound on the argument.
\begin{figure}[tbp]
\centering
\ViewHeading{(a) Geometric inclusions and compatible projections}
\begin{tikzpicture}
\node[vbox] (small) at (-4.2,0) {$\mathcal M^r$};
\node[vbox] (large) at (4.2,0) {$\mathcal M^R$};
\node[vresult] (base) at (0,-1.6) {$\mathcal M$};
\draw[varrow] (small)--(large) node[midway,above] {$j_{r,R}$};
\draw[varrow] (small)--(base) node[midway,below left] {$p_r$};
\draw[varrow] (large)--(base) node[midway,below right] {$p_R$};
\node[font=\footnotesize] at (0,-2.35) {$p_Rj_{r,R}=p_r$,\qquad $\mathcal L_{x,r}=p_r^*\mathcal L_x$};
\end{tikzpicture}
\ViewHeading{(b) Cohomology goes in the reverse direction}
\begin{tikzpicture}
\node[vbox,text width=4.9cm] (smallco) at (-3.8,0)
 {$H^1(\mathcal M^r,\{a\};\mathcal L_{x,r})$\\$\delta_{x,r}(1)$};
\node[vbox,text width=4.9cm] (largeco) at (3.8,0)
 {$H^1(\mathcal M^R,\{a\};\mathcal L_{x,R})$\\$\delta_{x,R}(1)$};
\draw[varrow] (largeco)--(smallco) node[midway,above] {$j_{r,R}^*\cong$};
\draw[viewteal,very thick] (-5.6,-1.7)--(5.6,-1.7);
\fill[viewteal] (-5.6,-1.7) circle (2.3pt);
\draw[viewteal,fill=white] (5.6,-1.7) circle (2.8pt);
\node[below] at (-5.6,-1.7) {$0$};
\node[below] at (5.6,-1.7) {$\lambda_-$};
\node[above,font=\footnotesize] at (0,-1.7)
 {fixed $x\neq0$: nonzero class throughout the guaranteed interval};
\node[font=\footnotesize,align=center] at (0,-2.85)
 {The open endpoint marks the limit of this projection argument,\\not a claimed death time for every obstruction.};
\end{tikzpicture}
\caption{Sub-reach transport holds for $0\leq r\leq R<\lambda_-$ and one
fixed character. Here $\mathcal M=\mathcal M_{\boldsymbol\lambda}$ and
$a=a_{\boldsymbol\lambda}$. Carrier points label local systems on the
whole carrier; they are not positions of separate holes in a single
system. Arithmetic powers are not the maps in this diagram.}
\label{fig:arithmetic-offset}
\end{figure}

Let
\[
 E_{\alpha,C}=R_{\alpha,C}\setminus\{0\},\qquad
 Z_{\alpha,C,\boldsymbol\lambda}
     =\iota_{\boldsymbol\lambda}(E_{\alpha,C}).
\]
The locus has no $r$ subscript: its definition does not change with the
offset radius.  Radius indexes transport of its associated classes.

\begin{theorem}[Selected spectrum on the carrier]\label{thm:dense-null}
For $\alpha>1/d$, the set $Z_{\alpha,C,\boldsymbol\lambda}$ is dense in
$\mathcal M_{\boldsymbol\lambda}$, is $\mathcal H^d$-null, and satisfies
\[
 \dim_{\mathrm H}Z_{\alpha,C,\boldsymbol\lambda}
          =\frac{d+1}{1+\alpha},\qquad
 \underline{\dim}_{\mathrm B}Z_{\alpha,C,\boldsymbol\lambda}
 =\overline{\dim}_{\mathrm B}Z_{\alpha,C,\boldsymbol\lambda}=d.
\]
Every selected parameter has the nonzero fixed-character obstruction of
\cref{prop:offset} at every sub-reach offset.
\end{theorem}

\begin{proof}
If $x$ is a rational character of order $q_0$, then
$b_{kq_0}(x)=0$ for all $k\geq1$.  Thus every nonzero rational character
belongs to $E_{\alpha,C}$, proving density.  Removing one point does not
alter the positive Hausdorff dimension in \cref{thm:dimension}.
The bi-Lipschitz bounds of \cref{prop:reach} transfer dimension and imply
nullity on the carrier.

For a bounded set $S$, let $N_s(S)$ be the least number of closed ambient
Euclidean balls of radius $s$ covering it.  Every finite closed-ball cover
of a dense subset of the compact carrier also covers its closure.  Hence
$N_s(Z_{\alpha,C,\boldsymbol\lambda})=N_s(\mathcal M_{\boldsymbol\lambda})$.
Grid covers and separated grids in the flat torus, transferred by
\cref{prop:reach}, give $N_s(\mathcal M_{\boldsymbol\lambda})\asymp s^{-d}$.
Both box dimensions therefore equal $d$.  The final claim is
\cref{prop:offset}.
\end{proof}

\begin{corollary}[Any prescribed positive codimension below $d$]
\label{cor:calibration}
For $0<\eta<d$, set $\alpha_{d,\eta}=(1+\eta)/(d-\eta)$.  Then
\[
 \alpha_{d,\eta}-\frac1d=\frac{\eta(d+1)}{d(d-\eta)}>0,\qquad
 \dim_{\mathrm H}Z_{\alpha_{d,\eta},C,\boldsymbol\lambda}=d-\eta.
\]
For $d=2$ and $\eta=1/(64\pi^2)$, this is
$2-1/(64\pi^2)=1.998416856505588\ldots$.
\end{corollary}

\begin{proof}
The first equality is direct algebra, and
$1+\alpha_{d,\eta}=(d+1)/(d-\eta)$.  Apply \cref{thm:dense-null}.
\end{proof}

The value $\eta$ is an input chosen by calibration, not a derived
universal constant.  It is a Hausdorff codimension, whereas $d-\eta$ is
the dimension.  Neither selects the radii $\boldsymbol\lambda$ or an
offset radius $r$.  The full nonextension parameter set is
$\T^d\setminus\{0\}$; $E_{\alpha,C}$ is a recurrence-selected subset,
not the entire locus of nonextension.  Carrier points label local systems
on the whole carrier; they are not spatial locations of different holes
inside one fixed local system.

\Cref{fig:arithmetic-dimensions} contrasts Hausdorff dimension, box
dimension, codimension, and the separate geometric radii.
\begin{figure}[tbp]
\centering
\begin{minipage}[t]{.48\linewidth}
\centering\ViewHeading{(a) Exact dimension formula for $d=2$}
\begin{tikzpicture}
\begin{axis}[viewaxis,width=7cm,height=5cm,xmin=0,xmax=4,ymin=0,ymax=2.25,
 xlabel={$\alpha$},ylabel={$h(\alpha)$},xtick={.5,1,2,3,4},ytick={0,1,2}]
\path[fill=viewgray!10] (axis cs:0,0) rectangle (axis cs:.5,2.25);
\addplot[viewblue,domain=.501:4,samples=70] {3/(1+x)};
\addplot[only marks,mark=o,viewblue] coordinates {(.5,2)};
\draw[dashed,viewgray] (axis cs:.5,0)--(axis cs:.5,2.25);
\node[font=\scriptsize,rotate=90] at (axis cs:.23,1.1) {outside stated theorem};
\node[font=\footnotesize,viewblue] at (axis cs:2.4,1.8) {$h(\alpha)=3/(1+\alpha)$};
\end{axis}
\end{tikzpicture}
\end{minipage}\hfill
\begin{minipage}[t]{.48\linewidth}
\centering\ViewHeading{(b) Dense and null are compatible}
\begin{tikzpicture}
\node[vbox,text width=5.9cm] at (0,0)
 {$Z=\iota_{\boldsymbol\lambda}(R_{\alpha,C}\setminus\{0\})$\\
 $\overline Z=\mathcal M_{\boldsymbol\lambda}$,\quad $\mathcal H^d(Z)=0$};
\node[vresult,text width=5.9cm] at (0,-1.5)
 {Hausdorff dimension: $h=(d+1)/(1+\alpha)$\\Both box dimensions: $d$};
\node[vinput,text width=5.9cm] at (0,-3.05)
 {Codimension $\eta=d-h$ is not a width.\\Radii $\lambda_j$ and offset $r$ are separate inputs.};
\end{tikzpicture}
\end{minipage}
\medskip
\begin{tabular}{@{}llll@{}}
\toprule
Quantity & $b_q$ & $\Fill_{q,x}(1)$ & $\widehat{\Fill}_{q,x}(1)$\\\midrule
Threshold & $<Cq^{-\alpha}$ & $>C^{-1}q^\alpha$ & $>C^{-1}q^{\alpha+1}$\\\bottomrule
\end{tabular}
\par\smallskip{\small Normalized component: $\|q^{-1}v_q\|_2<Cq^{-(\alpha+1)}$.}
\caption{Dimension, density, and threshold normalization. The curve is
the stated theorem's exact formula, not a fitted numerical estimate.
Dense subsets and their closures have the same finite closed-ball
covering numbers here, while their Hausdorff dimensions differ. No
positive-thickness fractal region is implied.}
\label{fig:arithmetic-dimensions}
\end{figure}

\section{Marked comparison and the limits of portability}
\label{sec:comparison}

Let $(K,a,\vartheta)$ be a connected finite CW complex with vertex $a$
and a specified epimorphism $\vartheta:\pi_1(K,a)\twoheadrightarrow\Z^d$.
Use a marked fibre $\C$ to define $\mathcal L_x^\vartheta$ by holonomy
$\rho_x\circ\vartheta$.

\begin{proposition}[Algebraic portability]\label{prop:portability}
If $x\ne0$, then
$H^0(K;\mathcal L_x^\vartheta)=0$ and
$\delta_x^\vartheta(1)\ne0$ in
$H^1(K,\{a\};\mathcal L_x^\vartheta)$.
\end{proposition}

\begin{proof}
Surjectivity makes $\rho_x\circ\vartheta$ nontrivial.  A global section
is a scalar fixed by every holonomy; one nonidentity scalar holonomy
forces it to vanish.  Since the vertex has $H^0=\C$, exactness of the
relative sequence makes its connecting homomorphism injective.
\end{proof}

\begin{corollary}[Rational tensor-power trivialization]
\label{cor:rational-power}
If $x\in\T^d$ satisfies $q_0x=0$ for an integer $q_0\geq1$, then
\[
 (\mathcal L_x^\vartheta)^{\otimes q_0}
       \cong\mathcal L_{q_0x}^\vartheta
       =\mathcal L_0^\vartheta,
 \qquad \delta_{q_0x}^\vartheta(1)=0.
\]
The isomorphism uses the marked fibre identification
$1^{\otimes q_0}\mapsto1$.  If $x\ne0$, the original class
$\delta_x^\vartheta(1)$ is nevertheless nonzero.
\end{corollary}

\begin{proof}
For every loop $g$, tensoring raises its scalar holonomy to the
$q_0$th power:
$\rho_x(\vartheta g)^{q_0}=\rho_{q_0x}(\vartheta g)=1$.
Thus the powered vertex state extends as the constant section, and its
connecting class vanishes.  The last assertion is \cref{prop:portability}.
\end{proof}

\begin{proposition}[Marked naturality]\label{prop:naturality}
Suppose a based map $f:(K',a')\to(K,a)$ and
$A\in\operatorname{GL}_d(\Z)$ satisfy
$\vartheta'=A\vartheta f_*$.  Put $\sigma_A(x)=A^{-T}x$.  Then, with
the marked fibre identifications,
\[
 f^*\mathcal L_x^\vartheta\cong
       \mathcal L_{\sigma_A(x)}^{\vartheta'},\qquad
 f^*\delta_x^\vartheta(1)=\delta_{\sigma_A(x)}^{\vartheta'}(1).
\]
A based homotopy equivalence of pairs induces an isomorphism on these
relative cohomology groups.
\end{proposition}

\begin{proof}
For $n\in\pi_1(K',a')$, the equality
$\langle A^{-T}x,A\vartheta(f_*n)\rangle
 =\langle x,\vartheta(f_*n)\rangle$ identifies the holonomy
representations.  Naturality of the relative long exact sequence proves
the connecting-class identity, and pair homotopy invariance proves the
last statement.
\end{proof}

Membership in the marked constant-$c$ approximation set transforms by
inclusions, not by an automatic equality:
\[
 \sigma_A(W_{\alpha,c})
     \subseteq W_{\alpha,\norm{A^{-T}}_{\infty\to\infty}c}.
\]
Applying the inverse transformation gives the corresponding inverse
bound.  The lattice is preserved, but the fixed norm and tolerance are
not generally preserved.  Likewise, \cref{prop:naturality} does not assert
that $\sigma_A$ is an isometry of boundary signatures or that the exact
constant-$C$ locus $R_{\alpha,C}$ is invariant.

Only the \emph{algebraic} extension obstruction ports to general $K$ in
\cref{prop:portability}.  No product-torus embedding, Euclidean chain
norm, single-row boundary formula, power-map diagram, or reach estimate
is supplied by an epimorphism of fundamental groups alone.  Those are
separate inputs or separate theorems.

\section{Related frameworks and further questions}
\label{sec:related}

\paragraph{Preorders and periodic spaces.}
Bubenik--de Silva--Scott \cite{bubenik-desilva-scott} formulate
interleavings for functors over preorders.  This provides a categorical
context for \cref{prop:feasible}, but we have not constructed an
interleaving metric or proved stability for our feasible-set functors.
Onus--Skraba \cite{onus-skraba} use bisheaves and persistent local systems
to study cycles in finite translation quotients of periodic complexes
and their lifting to the periodic space.  Our parameter $q$ instead
samples powers of a character on a fixed wedge.  A comparison with their
lifting problem would require additional maps and is not used in our
dimension proof.

\paragraph{Exact rank data and quantitative information.}
For example, Budur--Wang \cite{budur-wang} study cohomology jump loci of
rank-one local systems on smooth quasi-projective varieties.  Such loci
record cohomology-dimension conditions; our Boolean test records just
solvability of one marked boundary equation.  \Cref{prop:boolean}
concerns this particular sequence of tests, not all possible jump-locus
invariants or the structural theorems for algebraic varieties.

\paragraph{Precision and coefficient-dependent filling.}
Ghrist--Ding \cite{ghrist-ding} study network-sheaf cohomology filtered by
algebraic precision over $p$-adic integers and more generally discrete
valuation rings.  Our norm is Archimedean and our indexing combines
reverse divisibility with a real cost threshold.  Relating these two
settings would require a comparison of coefficient systems, filtrations,
and quantitative outputs, not merely the occurrence of holonomy in both.
Li--Manin \cite{li-manin} study homological filling functions of groups
and exhibit dependence of their asymptotic behavior on the coefficient
group.  Here the finite chain spaces, complex coefficients, selected
zero-chain, and Euclidean norm are fixed at each scale.  We do not
identify this reciprocal row norm with a group filling function.

\paragraph{Three further questions.}
The following are directions for additional work, not consequences of
the present theorems:
\begin{enumerate}
 \item Which changes of marking can be quotiented out while retaining
 a nondegenerate metric and controlled changes in the exact tolerance?
 The inverse-transpose comparison of \cref{prop:naturality} supplies
 algebraic transport, not an isometry of the chosen signatures.
 \item Can a finite cellular-inclusion construction reproduce the
 relevant filling profiles, with an explicit reindexing and quantitative
 comparison to the reverse-divisibility diagram?  Its non-inclusion maps
 and the threshold-direction constraint must be addressed.
 \item For larger finite CW complexes, which hypotheses on twisted
 boundary matrices and the marked cycle control minimum filling norms
 through their singular values, and yield a tractable recurrence locus?
 Neither a dimension formula nor an exact scalar-row reduction is assumed.
\end{enumerate}

\section{Scope and reproducibility}

The construction proves a marked, normed, finite-at-each-scale arithmetic
model.  It does not prove an invariant of arbitrary unmarked homotopy
types, a general correspondence with inclusion filtrations, a physical
channel law, or a mechanism selecting $1/(64\pi^2)$.  A change of
coefficient convention is handled explicitly by dualization; a change of
norm, basis, or arithmetic indexing requires a new comparison argument.

The three forms of transport in the series remain distinct: arithmetic
maps are chain maps between different scales, offset maps give
contravariant relative-cohomology isomorphisms for fixed characters, and
finite-response marks are endomorphisms of a supplied fixed quotient.
No automatic conversion between them is used.

The shared script \path{scripts/verify_restructured_series.py} checks exact
polynomial identities, rational torsion cases, finite-field examples,
and sampled numerical filling, signature, and carrier identities.  It
also checks the two-dimensional powered-torsion example and finite exact
instances of the constructive arithmetic estimates in the appendices.  Its
finite checks do not certify an infinitely-often condition, a Hausdorff
dimension formula, or a theorem for all parameters.  Those assertions
rest on the proofs and the stated classical dimension theorem.

The deterministic tables in \cref{fig:arithmetic-profiles,fig:arithmetic-signature}
are included in \path{figures/arithmetic_plot_tables.tex}; a normal LaTeX
build requires no external data files or Python execution. They can be
regenerated or checked by \path{scripts/build_series_plot_data.py}.
The figures use native TikZ and PGFPlots; their finite samples do not
strengthen the scope of the theorems.

\appendix
\section{Normalization and the earlier constant-one locus}
\label{sec:legacy}

Keep the torus metric, common denominator, and unreduced-denominator
convention of \cref{sec:overview,sec:dimension}.  The earlier cube-based
set was
\[
 W_d^{\mathrm{cube}}(\tau)
 =\left\{\xi\in[0,1]^d:
 \norm{\xi-p/q}_\infty<q^{-(1+\tau)}
 \text{ for infinitely many }(p,q)\in\Z^d\times\N\right\}.
\]
If $\pi:\R^d\to\T^d$ is the quotient map, then
\begin{equation}\label{eq:legacy-locus}
 \pi(W_d^{\mathrm{cube}}(\alpha))=W_{\alpha,1},\qquad
 E^{\mathrm{sup}}_\alpha=W_{\alpha,1}\setminus\{0\}.
\end{equation}
Indeed, multiplying the cube inequality by $q$ gives exactly the
definition of $\Delta_q$; conversely each torus parameter has a lift in
the cube and the inequality is independent of that lift.  For each fixed
$q$ only finitely many $p$ can satisfy the cube inequality, so infinitely
many pairs mean infinitely many denominators.

The removal of $0$ in \cref{eq:legacy-locus} occurs \emph{after} taking
the quotient.  Removing only the cube point $(0,\ldots,0)$ is insufficient:
every corner in $\{0,1\}^d$ represents the same trivial character.
Equivalently, remove all those corners before taking the periodic image.

The old Toric exponent was $\tau_{\mathrm{Toric}}=\alpha$, while the
old Quantitative exponent was $\tau_{\mathrm{Quantitative}}=\alpha-1$.
The old sup-norm recurrence test is exactly $\Delta_q<q^{-\alpha}$;
the present chord-norm test is $b_q<Cq^{-\alpha}$.  The comparison is
\cref{eq:locus-sandwich}, not equality at a common constant.  Equal
Hausdorff dimensions therefore do not identify the selected sets.

\begin{proposition}[Rational membership and density for all positive exponents]
\label{prop:legacy-density}
For every $\alpha,c,C>0$, both $W_{\alpha,c}$ and $R_{\alpha,C}$
contain every rational character.  Removing $0$ leaves a dense subset
of $\T^d$ in either case.
\end{proposition}

\begin{proof}
For a rational character choose $q_0$ with $q_0x=0$.  Then
$\Delta_{kq_0}(x)=b_{kq_0}(x)=0$ for every $k\geq1$, satisfying both
strict inequalities with positive tolerances.  Nonzero rational
characters are dense in the torus, which proves the final assertion.
\end{proof}

\begin{corollary}[The earlier sup-norm spectrum on the same carrier]
\label{cor:legacy-carrier}
Let $\alpha>1/d$ and $\boldsymbol\lambda>0$.  The set
$Z^{\mathrm{sup}}_{\alpha,\boldsymbol\lambda}
 =\iota_{\boldsymbol\lambda}(E^{\mathrm{sup}}_\alpha)$ is dense in
$\mathcal M_{\boldsymbol\lambda}$, is $\mathcal H^d$-null, and has
\[
 \dim_{\mathrm H}Z^{\mathrm{sup}}_{\alpha,\boldsymbol\lambda}
   =\frac{d+1}{1+\alpha},\qquad
 \underline{\dim}_{\mathrm B}Z^{\mathrm{sup}}_{\alpha,\boldsymbol\lambda}
 =\overline{\dim}_{\mathrm B}Z^{\mathrm{sup}}_{\alpha,\boldsymbol\lambda}=d.
\]
Every parameter in this set labels a nonzero fixed-character connecting
class at every offset $0\leq r<\lambda_-$.
\end{corollary}

\begin{proof}
\Cref{lem:constants} gives the positive dimension of $W_{\alpha,1}$;
removing one point does not change it.  The bi-Lipschitz map of
\cref{prop:reach} transfers that dimension, whose strict inequality with
$d$ gives nullity.  Density follows from \cref{prop:legacy-density}.
As in the proof of \cref{thm:dense-null}, a finite closed-ball cover of
a dense subset covers the compact carrier, whose covering numbers are
comparable to $s^{-d}$.  This gives both box dimensions.  Finally apply
\cref{prop:offset} to each nonzero character.
\end{proof}

For $\alpha=(1+\eta)/(d-\eta)$, $0<\eta<d$, this restores the earlier
constant-one spectrum of dimension $d-\eta$, including the stated
two-dimensional calibration.  Either $E^{\mathrm{sup}}_\alpha$ or
$R_{\alpha,C}\setminus\{0\}$ can serve as the precisely specified
target of the conditional bridge in the foundations companion; a
bi-Lipschitz source map and coefficient/seed identifications remain
additional inputs.  The all-positive-exponent proposition above makes no
endpoint nullity or endpoint Boolean no-go claim.

\section{Estimates for the explicit arithmetic comparison}
\label{sec:explicit-estimates}

We prove the claims of \cref{ex:explicit-arithmetic} without appealing
to a dimension theorem.  Put
\[
 \ell_n=\sum_{m=1}^n10^{-m!},\qquad q_n=10^{n!}.
\]
Then $q_n\ell_n\in\Z$.  Since the factorial exponents in the tail
increase by at least one, comparison with a geometric series yields
\[
 0<\ell-\ell_n<2\,10^{-(n+1)!},\qquad
 \dist(q_n\ell,\Z)<2q_n^{-n}.
\]
The number $\ell$ is irrational: if $\ell=a/b$, its positive difference
from $\ell_n$ would be at least $1/(bq_n)$, contradicting this tail
bound for sufficiently large $n$.  Thus $x_{\mathrm L}$ is non-torsion.
The upper chord bound gives
$b_{q_n}(x_{\mathrm L})<4\pi q_n^{-n}$, which is less than
$Cq_n^{-\alpha}$ for all sufficiently large $n$, for any fixed
$\alpha,C>0$.

For the other parameter, let $p$ be a nearest integer to $q\gamma$ and
write $\gamma'=(-1-\sqrt5)/2$.  The integer $p^2+pq-q^2$ is nonzero
because the roots of $t^2+t-1$ are irrational.  Therefore
\[
 \left|\frac pq-\gamma\right|
 \left|\frac pq-\gamma'\right|
 =\frac{|p^2+pq-q^2|}{q^2}\geq\frac1{q^2}.
\]
The second factor is at most $\sqrt5+1/(2q)<3$, so
$\dist(q\gamma,\Z)>1/(3q)$.  The lower chord bound now gives
\[
 b_q(x_{\mathrm G})>\frac4{3q}.
\]
For $\alpha>1$ this is eventually greater than $Cq^{-\alpha}$;
hence $x_{\mathrm G}\notin R_{\alpha,C}$.  Irrationality of $\gamma$
makes its Boolean sequence identical to that of $x_{\mathrm L}$.
This estimate does not exclude recurrence when $\alpha\leq1$ and is
not used for that purpose.

\clearpage
\bibliographystyle{plain}
\bibliography{series_references}
\end{document}